\chapter*{ESPECIFICACIONES DE MATERIALES}
\addcontentsline{toc}{chapter}{ESPECIFICACIONES DE MATERIALES}

% Reinicio de contadores para numeracion coherente con el capitulo 3.
\setcounter{section}{0}
\renewcommand{\thesection}{3.\arabic{section}.}
\renewcommand{\thesubsection}{3.\arabic{section}.\arabic{subsection}.}
\renewcommand{\thesubsubsection}{3.\arabic{section}.\arabic{subsection}.\arabic{subsubsection}.}
\renewcommand{\theHsection}{materiales.\arabic{section}}
\renewcommand{\theHsubsection}{materiales.\arabic{section}.\arabic{subsection}}
\renewcommand{\theHsubsubsection}{materiales.\arabic{section}.\arabic{subsection}.\arabic{subsubsection}}

\section{CONDUCTORES ELECTRICOS}

\subsection{CABLES DE ALIMENTACION (TIPO N2XOH o NH-80)}
Los conductores del alimentador principal y de los subalimentadores (si los hubiera) seran de cobre electrolitico de temple blando, cableado concentrico y con aislamiento termoplastico libre de halogenos.
\begin{itemize}
  \item \textbf{Material:} Cobre electrolitico de temple blando (99.9\% de pureza).
  \item \textbf{Aislamiento:} Compuesto termoestable libre de halogenos y de baja emision de humos toxicos y corrosivos (XLPE / LSOH).
  \item \textbf{Normas de fabricacion:} NTP-IEC 60228 (conductores), NTP-IEC 60502-1, IEC 60754-1, IEC 60754-2, IEC 61034-2 e IEC 60332-3-24 (no propagacion del incendio).
  \item \textbf{Temperatura de operacion:} Hasta \(90^\circ\text{C}\) en servicio continuo, \(130^\circ\text{C}\) en sobrecarga y \(250^\circ\text{C}\) en cortocircuito.
\end{itemize}

\subsection{CONDUCTORES PARA CIRCUITOS DERIVADOS (TIPO THW o NH-80)}
Los cables para el cableado de circuitos de alumbrado y tomacorrientes seran unipolares de cobre, con aislamiento de PVC retardante de la llama.
\begin{itemize}
  \item \textbf{Aislamiento:} Cloruro de polivinilo (PVC) resistente al calor y a la humedad.
  \item \textbf{Norma de fabricacion:} NTP 370.252 / NTP 370.053.
  \item \textbf{Temperatura de operacion:} Hasta \(75^\circ\text{C}\) en ambientes secos y humedos.
  \item \textbf{Secciones minimas:} Alumbrado (\(2.5\text{ mm}^2\)), tomacorrientes (\(2.5\text{ mm}^2\)), cocina (\(10.0\text{ mm}^2\)), lavadora (\(10.0\text{ mm}^2\)).
\end{itemize}

\section{CANALIZACIONES Y ACCESORIOS}

\subsection{TUBERIAS DE PVC CLASE PESADA (PVC-P)}
Las tuberias para los circuitos empotrados seran rigidas de Cloruro de Polivilino (PVC) no plastificado, clase pesada.
\begin{itemize}
  \item \textbf{Caracteristicas:} Resistentes al impacto, aplastamiento, humedad y agentes quimicos. Retardantes de la llama y autoextinguibles.
  \item \textbf{Norma de fabricacion:} NTP 399.006 (INDECOPI ex ITINTEC).
  \item \textbf{Longitud y union:} Tubos de 3.00 m de longitud con campana en un extremo para union mediante pegamento PVC.
  \item \textbf{Diametros minimos:} Derivados (\(20\text{ mm}\) o 3/4''), alimentador general (\(40\text{ mm}\) o 1 1/2''), subalimentador a T2 (\(32\text{ mm}\) o 1 1/4'').
\end{itemize}

\subsection{CAJAS METALICAS DE SALIDA Y DISPOSITIVOS}
Las cajas para tomacorrientes, interruptores y centros de luz seran de fierro galvanizado de construccion embutida en una sola pieza.
\begin{itemize}
  \item \textbf{Espesor de plancha:} No menor a \(1.6\text{ mm}\).
  \item \textbf{Galvanizacion:} Galvanizado en caliente, segun ASTM A526-71 designacion G-90, con un recubrimiento minimo de 40\% de zinc.
  \item \textbf{Formas y usos:}
    - \textbf{Rectangulares (100x55x50 mm):} Para interruptores y tomacorrientes de pared.
    - \textbf{Octogonales (100x50 mm):} Para centros de iluminacion en techo o pared.
\end{itemize}

\subsection{CAJAS DE PASO Y ESTANCAS}
Fabricadas en plancha de fierro galvanizado de un espesor no menor a \(1.8\text{ mm}\) para cajas con dimensiones mayores a 200 mm. Para exteriores (bomba de agua) se utilizaran cajas estancas termoplasticas con grado de proteccion minimo IP55.

\section{DISPOSITIVOS DE MANIOBRA Y PROTECCION}

\subsection{INTERRUPTORES TERMOMAGNETICOS RESIDENCIALES}
Dispositivos automaticos termicos y magneticos destinados a la proteccion contra sobrecargas y cortocircuitos.
\begin{itemize}
  \item \textbf{Tipo de montaje:} Riel DIN de 35 mm, de operacion manual y disparo libre.
  \item \textbf{Capacidad de ruptura:} Minimo \(10\text{ kA}\) a 220 V (NTP-IEC 60898-1).
  \item \textbf{Curva de disparo:} Curva tipo C (apto para cargas residenciales con corrientes de arranque bajas).
  \item \textbf{Marcado:} Indicar marca, corriente nominal (A), tension (V) e indicador visual de estado ON/OFF.
\end{itemize}

\subsection{INTERRUPTORES DIFERENCIALES}
Interruptores diferenciales para la proteccion contra fugas de corriente a tierra y proteccion de personas contra contactos directos e indirectos.
\begin{itemize}
  \item \textbf{Sensibilidad residual:} \(30\text{ mA}\).
  \item \textbf{Tiempo de respuesta:} Menor a \(30\text{ ms}\) a la corriente nominal de defecto.
  \item \textbf{Capacidad de corriente nominal:} Mayor o igual a la corriente nominal del termomagnetico asociado (NTP-IEC 61008-1).
\end{itemize}

\subsection{TOMACORRIENTES Y ACCESORIOS}
Los tomacorrientes de pared seran bipolares dobles con contacto de tierra lateral (Schuko) o universal de tres polos, con alveolos protegidos contra contactos accidentales.
\begin{itemize}
  \item \textbf{Capacidad y Tension:} \(16\text{ A}\) a 250 V AC, 60 Hz. Para salidas especiales (cocina, lavadora) se emplearan tomacorrientes de mayor capacidad (20 A o 32 A).
  \item \textbf{Placa protectora:} Placa de policarbonato autoextinguible de alta resistencia.
\end{itemize}

\section{SISTEMA DE PUESTA A TIERRA}

\subsection{ELECTRODO VERTICAL}
Sera una varilla de acero recubierta de cobre por proceso electrolitico (tipo copperweld), con un diametro minimo de \(19\text{ mm}\) (\(3/4''\)) y una longitud de \(2.40\text{ m}\).

\subsection{CONDUCTOR DE CONEXION Y ACCESORIOS}
\begin{itemize}
  \item \textbf{Conductor principal:} Cobre electrolitico desnudo de temple medio de \(10\text{ mm}^2\) de seccion.
  \item \textbf{Caja de registro:} De concreto prefabricado con tapa de registro (0.30x0.30 m) instalada en piso terminado para inspeccion.
  \item \textbf{Conector de varilla:} Abrazadera de bronce tipo perno en U o conector certificado para asegurar contacto franco de baja resistencia.
  \item \textbf{Tratamiento del terreno:} Aplicacion de dosis quimica de gel o sales conductivas no corrosivas para mejorar la resistividad del terreno.
\end{itemize}

\section{ACCESORIOS COMPLEMENTARIOS}

\subsection{CINTA AISLANTE DE PVC}
Cinta autoadhesiva elastoplastica de cloruro de polivinilo, de 20 mm de ancho, 10 m de largo y 0.5 mm de espesor minimo, con una rigidez dielectrica minima de 13.8 kV/mm y apta para operar a \(80^\circ\text{C}\).

\subsection{LUMINARIAS LED}
Las luminarias seran paneles LED redondos de adosar de 18 W de potencia nominal, con una eficiencia luminosa minima de 90 lm/W y vida util proyectada de 25,000 horas.
